\IfFileExists{stacks-project.cls}{%
\documentclass{stacks-project}
}{%
\documentclass{amsart}
}

% For dealing with references we use the comment environment
\usepackage{verbatim}
\newenvironment{reference}{\comment}{\endcomment}
%\newenvironment{reference}{}{}
\newenvironment{slogan}{\comment}{\endcomment}
\newenvironment{history}{\comment}{\endcomment}

% For commutative diagrams we use Xy-pic
\usepackage[all]{xy}

% We use 2cell for 2-commutative diagrams.
\xyoption{2cell}
\UseAllTwocells

% We use multicol for the list of chapters between chapters
\usepackage{multicol}

% This is generally recommended for better output
\usepackage{lmodern}
\usepackage[T1]{fontenc}

% For cross-file-references
\usepackage{xr-hyper}

% Package for hypertext links:
\usepackage{hyperref}

% For any local file, say "hello.tex" you want to link to please
% use \externaldocument[hello-]{hello}
\externaldocument[introduction-]{introduction}
\externaldocument[conventions-]{conventions}
\externaldocument[sets-]{sets}
\externaldocument[categories-]{categories}
\externaldocument[topology-]{topology}
\externaldocument[sheaves-]{sheaves}
\externaldocument[sites-]{sites}
\externaldocument[stacks-]{stacks}
\externaldocument[fields-]{fields}
\externaldocument[algebra-]{algebra}
\externaldocument[brauer-]{brauer}
\externaldocument[homology-]{homology}
\externaldocument[derived-]{derived}
\externaldocument[simplicial-]{simplicial}
\externaldocument[more-algebra-]{more-algebra}
\externaldocument[smoothing-]{smoothing}
\externaldocument[modules-]{modules}
\externaldocument[sites-modules-]{sites-modules}
\externaldocument[injectives-]{injectives}
\externaldocument[cohomology-]{cohomology}
\externaldocument[sites-cohomology-]{sites-cohomology}
\externaldocument[dga-]{dga}
\externaldocument[dpa-]{dpa}
\externaldocument[sdga-]{sdga}
\externaldocument[hypercovering-]{hypercovering}
\externaldocument[schemes-]{schemes}
\externaldocument[constructions-]{constructions}
\externaldocument[properties-]{properties}
\externaldocument[morphisms-]{morphisms}
\externaldocument[coherent-]{coherent}
\externaldocument[divisors-]{divisors}
\externaldocument[limits-]{limits}
\externaldocument[varieties-]{varieties}
\externaldocument[topologies-]{topologies}
\externaldocument[descent-]{descent}
\externaldocument[perfect-]{perfect}
\externaldocument[more-morphisms-]{more-morphisms}
\externaldocument[flat-]{flat}
\externaldocument[groupoids-]{groupoids}
\externaldocument[more-groupoids-]{more-groupoids}
\externaldocument[etale-]{etale}
\externaldocument[chow-]{chow}
\externaldocument[intersection-]{intersection}
\externaldocument[pic-]{pic}
\externaldocument[weil-]{weil}
\externaldocument[adequate-]{adequate}
\externaldocument[dualizing-]{dualizing}
\externaldocument[duality-]{duality}
\externaldocument[discriminant-]{discriminant}
\externaldocument[derham-]{derham}
\externaldocument[local-cohomology-]{local-cohomology}
\externaldocument[algebraization-]{algebraization}
\externaldocument[curves-]{curves}
\externaldocument[resolve-]{resolve}
\externaldocument[models-]{models}
\externaldocument[functors-]{functors}
\externaldocument[equiv-]{equiv}
\externaldocument[pione-]{pione}
\externaldocument[etale-cohomology-]{etale-cohomology}
\externaldocument[proetale-]{proetale}
\externaldocument[relative-cycles-]{relative-cycles}
\externaldocument[more-etale-]{more-etale}
\externaldocument[trace-]{trace}
\externaldocument[crystalline-]{crystalline}
\externaldocument[spaces-]{spaces}
\externaldocument[spaces-properties-]{spaces-properties}
\externaldocument[spaces-morphisms-]{spaces-morphisms}
\externaldocument[decent-spaces-]{decent-spaces}
\externaldocument[spaces-cohomology-]{spaces-cohomology}
\externaldocument[spaces-limits-]{spaces-limits}
\externaldocument[spaces-divisors-]{spaces-divisors}
\externaldocument[spaces-over-fields-]{spaces-over-fields}
\externaldocument[spaces-topologies-]{spaces-topologies}
\externaldocument[spaces-descent-]{spaces-descent}
\externaldocument[spaces-perfect-]{spaces-perfect}
\externaldocument[spaces-more-morphisms-]{spaces-more-morphisms}
\externaldocument[spaces-flat-]{spaces-flat}
\externaldocument[spaces-groupoids-]{spaces-groupoids}
\externaldocument[spaces-more-groupoids-]{spaces-more-groupoids}
\externaldocument[bootstrap-]{bootstrap}
\externaldocument[spaces-pushouts-]{spaces-pushouts}
\externaldocument[spaces-chow-]{spaces-chow}
\externaldocument[groupoids-quotients-]{groupoids-quotients}
\externaldocument[spaces-more-cohomology-]{spaces-more-cohomology}
\externaldocument[spaces-simplicial-]{spaces-simplicial}
\externaldocument[spaces-duality-]{spaces-duality}
\externaldocument[formal-spaces-]{formal-spaces}
\externaldocument[restricted-]{restricted}
\externaldocument[spaces-resolve-]{spaces-resolve}
\externaldocument[formal-defos-]{formal-defos}
\externaldocument[defos-]{defos}
\externaldocument[cotangent-]{cotangent}
\externaldocument[examples-defos-]{examples-defos}
\externaldocument[algebraic-]{algebraic}
\externaldocument[examples-stacks-]{examples-stacks}
\externaldocument[stacks-sheaves-]{stacks-sheaves}
\externaldocument[criteria-]{criteria}
\externaldocument[artin-]{artin}
\externaldocument[quot-]{quot}
\externaldocument[stacks-properties-]{stacks-properties}
\externaldocument[stacks-morphisms-]{stacks-morphisms}
\externaldocument[stacks-limits-]{stacks-limits}
\externaldocument[stacks-cohomology-]{stacks-cohomology}
\externaldocument[stacks-perfect-]{stacks-perfect}
\externaldocument[stacks-introduction-]{stacks-introduction}
\externaldocument[stacks-more-morphisms-]{stacks-more-morphisms}
\externaldocument[stacks-geometry-]{stacks-geometry}
\externaldocument[moduli-]{moduli}
\externaldocument[moduli-curves-]{moduli-curves}
\externaldocument[examples-]{examples}
\externaldocument[exercises-]{exercises}
\externaldocument[guide-]{guide}
\externaldocument[desirables-]{desirables}
\externaldocument[coding-]{coding}
\externaldocument[obsolete-]{obsolete}
\externaldocument[fdl-]{fdl}
\externaldocument[index-]{index}

% Theorem environments.
%
\theoremstyle{plain}
\newtheorem{theorem}[subsection]{Theorem}
\newtheorem{proposition}[subsection]{Proposition}
\newtheorem{lemma}[subsection]{Lemma}

\theoremstyle{definition}
\newtheorem{definition}[subsection]{Definition}
\newtheorem{example}[subsection]{Example}
\newtheorem{exercise}[subsection]{Exercise}
\newtheorem{situation}[subsection]{Situation}

\theoremstyle{remark}
\newtheorem{remark}[subsection]{Remark}
\newtheorem{remarks}[subsection]{Remarks}

\numberwithin{equation}{subsection}

% Macros
%
\def\lim{\mathop{\mathrm{lim}}\nolimits}
\def\colim{\mathop{\mathrm{colim}}\nolimits}
\def\Spec{\mathop{\mathrm{Spec}}}
\def\Hom{\mathop{\mathrm{Hom}}\nolimits}
\def\Ext{\mathop{\mathrm{Ext}}\nolimits}
\def\SheafHom{\mathop{\mathcal{H}\!\mathit{om}}\nolimits}
\def\SheafExt{\mathop{\mathcal{E}\!\mathit{xt}}\nolimits}
\def\Sch{\mathit{Sch}}
\def\Mor{\mathop{\mathrm{Mor}}\nolimits}
\def\Ob{\mathop{\mathrm{Ob}}\nolimits}
\def\Sh{\mathop{\mathit{Sh}}\nolimits}
\def\NL{\mathop{N\!L}\nolimits}
\def\CH{\mathop{\mathrm{CH}}\nolimits}
\def\proetale{{pro\text{-}\acute{e}tale}}
\def\etale{{\acute{e}tale}}
\def\QCoh{\mathit{QCoh}}
\def\Ker{\mathop{\mathrm{Ker}}}
\def\Im{\mathop{\mathrm{Im}}}
\def\Coker{\mathop{\mathrm{Coker}}}
\def\Coim{\mathop{\mathrm{Coim}}}

% Boxtimes
%
\DeclareMathSymbol{\boxtimes}{\mathbin}{AMSa}{"02}

%
% Macros for moduli stacks/spaces
%
\def\QCohstack{\mathcal{QC}\!\mathit{oh}}
\def\Cohstack{\mathcal{C}\!\mathit{oh}}
\def\Spacesstack{\mathcal{S}\!\mathit{paces}}
\def\Quotfunctor{\mathrm{Quot}}
\def\Hilbfunctor{\mathrm{Hilb}}
\def\Curvesstack{\mathcal{C}\!\mathit{urves}}
\def\Polarizedstack{\mathcal{P}\!\mathit{olarized}}
\def\Complexesstack{\mathcal{C}\!\mathit{omplexes}}
% \Pic is the operator that assigns to X its picard group, usage \Pic(X)
% \Picardstack_{X/B} denotes the Picard stack of X over B
% \Picardfunctor_{X/B} denotes the Picard functor of X over B
\def\Pic{\mathop{\mathrm{Pic}}\nolimits}
\def\Picardstack{\mathcal{P}\!\mathit{ic}}
\def\Picardfunctor{\mathrm{Pic}}
\def\Deformationcategory{\mathcal{D}\!\mathit{ef}}
\usepackage{cleveref}
\crefname{originalsection}{original section}{original sections}
\Crefname{originalsection}{Original section}{Original sections}
\makeatletter
\newlabel{bootstrap-section-representable-by-spaces-properties}{{046E}{source}{}{sourcepointer-1}{}}
\newlabel{bootstrap-section-representable-by-spaces-properties@cref}{{[originalsection][0][]046E}{[1][0][]source}{}{}{}}
\makeatother
\crefname{originalequation}{original equation}{original equations}
\Crefname{originalequation}{Original equation}{Original equations}
\makeatletter
\newlabel{equation-etale-topos}{{0DEB}{source}{}{sourcepointer-2}{}}
\newlabel{equation-etale-topos@cref}{{[originalequation][0][]0DEB}{[1][0][]source}{}{}{}}
\makeatother
\crefname{originalproposition}{original proposition}{original propositions}
\Crefname{originalproposition}{Original proposition}{Original propositions}
\makeatletter
\newlabel{etale-cohomology-proposition-quasi-coherent-sheaf-fpqc}{{03OG}{source}{}{sourcepointer-3}{}}
\newlabel{etale-cohomology-proposition-quasi-coherent-sheaf-fpqc@cref}{{[originalproposition][0][]03OG}{[1][0][]source}{}{}{}}
\makeatother
\crefname{originallemma}{original lemma}{original lemmas}
\Crefname{originallemma}{Original lemma}{Original lemmas}
\makeatletter
\newlabel{lemma-affine-representable-by-algebraic-spaces}{{0AKN}{source}{}{sourcepointer-4}{}}
\newlabel{lemma-affine-representable-by-algebraic-spaces@cref}{{[originallemma][0][]0AKN}{[1][0][]source}{}{}{}}
\makeatother
\crefname{originallemma}{original lemma}{original lemmas}
\Crefname{originallemma}{Original lemma}{Original lemmas}
\makeatletter
\newlabel{lemma-identify-affine-etale}{{0DEH}{source}{}{sourcepointer-5}{}}
\newlabel{lemma-identify-affine-etale@cref}{{[originallemma][0][]0DEH}{[1][0][]source}{}{}{}}
\makeatother
\crefname{originallemma}{original lemma}{original lemmas}
\Crefname{originallemma}{Original lemma}{Original lemmas}
\makeatletter
\newlabel{lemma-permanence-representable}{{0AQ0}{source}{}{sourcepointer-6}{}}
\newlabel{lemma-permanence-representable@cref}{{[originallemma][0][]0AQ0}{[1][0][]source}{}{}{}}
\makeatother
\crefname{originallemma}{original lemma}{original lemmas}
\Crefname{originallemma}{Original lemma}{Original lemmas}
\makeatletter
\newlabel{spaces-properties-lemma-etale-permanence}{{03FV}{source}{}{sourcepointer-7}{}}
\newlabel{spaces-properties-lemma-etale-permanence@cref}{{[originallemma][0][]03FV}{[1][0][]source}{}{}{}}
\makeatother
\crefname{originalremark}{original remark}{original remarks}
\Crefname{originalremark}{Original remark}{Original remarks}
\makeatletter
\newlabel{remark-set-theoretic}{{0AIS}{source}{}{sourcepointer-8}{}}
\newlabel{remark-set-theoretic@cref}{{[originalremark][0][]0AIS}{[1][0][]source}{}{}{}}
\makeatother
\title{Every formal algebraic space carries a structure sheaf of topological rings}
\author{The Stacks Project Authors}
\date{stacks-project:formal-spaces.tex:local-snapshot-20261012; GNU FDL 1.2 or later}
\begin{document}
\maketitle
\noindent\textit{Editorial note.} All delivered mathematical text is the author's own, copied from the pinned Stacks Project chapter file formal-spaces.tex; no statement or proof step was rewritten, reordered or removed.\par\medskip
\noindent\textit{Editorial note.} The proof displays three exact sequences of sections, two of them inside the limit computation; all lie inside the counted span and are reproduced verbatim.\par\medskip
\noindent\textit{Editorial note.} The last paragraph compares two fundamental systems of open submodules, one on the sections over the formal space and one on the product over the members of the covering; its final clause is reproduced exactly as delivered, without any rewording or correction.\par\medskip
\noindent\textit{Editorial note.} The counted spans carry 8 navigation stubs; no bibliography citation occurs in the spans.\par\medskip
\noindent\textit{Editorial note.} The preamble span points at the chapter's own preamble.tex shipped with the frozen source, used verbatim as staged.\par\medskip
\noindent\textit{Allowed standard prerequisites (editorial).}\begin{itemize}
\item The small etale site of a formal algebraic space and the identification with affine formal algebraic spaces etale over it. Same chapter, Equation equation-etale-topos at proof lines 7669-7671 and Lemma lemma-identify-affine-etale at proof lines 7674-7677. Reduces the construction to a site whose objects are concrete, and identifies reduction as an equivalence with that site.
\item Affine formal algebraic spaces as filtered colimits of their defining affine schemes. Same chapter, Definition definition-affine-formal-algebraic-space, invoked at proof lines 7694-7695. Supplies the filtered system over which the exact sequences are taken.
\item Permanence of representability and affineness of a representable morphism. Same chapter, Lemma lemma-permanence-representable at proof lines 7686-7688 and Lemma lemma-affine-representable-by-algebraic-spaces at proof lines 7688-7689. Shows that the members of a covering are representable and hence affine, so that their fibre products are affine schemes.
\item Etale morphisms are preserved by the passage considered, and the structure sheaf satisfies the sheaf condition for the etale topology. Properties of Spaces, Lemma spaces-properties-lemma-etale-permanence at proof lines 7689-7692 and Etale Cohomology, Proposition etale-cohomology-proposition-quasi-coherent-sheaf-fpqc at proof lines 7716-7720. Supplies the exact sequences whose filtered limit gives the sheaf condition at the level of the formal space.
\end{itemize}\medskip
\ifdefined\hypertarget\hypertarget{sourcepointer-1}{}\fi
\noindent\textit{Source reference.} Original section 046E: sources/stacks-chow/tags.txt, source lines 5469--5469. Pointer outside the counted span, invoked once for context; the load-bearing argument of the counted statement is the author's proof reproduced in full in the proof span.\par\medskip
\ifdefined\hypertarget\hypertarget{sourcepointer-2}{}\fi
\noindent\textit{Source reference.} Original equation 0DEB: sources/stacks-formal-spaces/pinned-original-source.tex, source lines 7413--7413. Pointer outside the counted span, invoked once for context; the load-bearing argument of the counted statement is the author's proof reproduced in full in the proof span.\par\medskip
\ifdefined\hypertarget\hypertarget{sourcepointer-3}{}\fi
\noindent\textit{Source reference.} Original proposition 03OG: sources/stacks-chow/tags.txt, source lines 4821--4821. Pointer outside the counted span, invoked once for context; the load-bearing argument of the counted statement is the author's proof reproduced in full in the proof span.\par\medskip
\ifdefined\hypertarget\hypertarget{sourcepointer-4}{}\fi
\noindent\textit{Source reference.} Original lemma 0AKN: sources/stacks-formal-spaces/pinned-original-source.tex, source lines 4209--4209. Pointer outside the counted span, invoked once for context; the load-bearing argument of the counted statement is the author's proof reproduced in full in the proof span.\par\medskip
\ifdefined\hypertarget\hypertarget{sourcepointer-5}{}\fi
\noindent\textit{Source reference.} Original lemma 0DEH: sources/stacks-formal-spaces/pinned-original-source.tex, source lines 7614--7614. Pointer outside the counted span, invoked once for context; the load-bearing argument of the counted statement is the author's proof reproduced in full in the proof span.\par\medskip
\ifdefined\hypertarget\hypertarget{sourcepointer-6}{}\fi
\noindent\textit{Source reference.} Original lemma 0AQ0: sources/stacks-formal-spaces/pinned-original-source.tex, source lines 4076--4076. Pointer outside the counted span, invoked once for context; the load-bearing argument of the counted statement is the author's proof reproduced in full in the proof span.\par\medskip
\ifdefined\hypertarget\hypertarget{sourcepointer-7}{}\fi
\noindent\textit{Source reference.} Original lemma 03FV: sources/stacks-chow/tags.txt, source lines 4506--4506. Pointer outside the counted span, invoked once for context; the load-bearing argument of the counted statement is the author's proof reproduced in full in the proof span.\par\medskip
\ifdefined\hypertarget\hypertarget{sourcepointer-8}{}\fi
\noindent\textit{Source reference.} Original remark 0AIS: sources/stacks-formal-spaces/pinned-original-source.tex, source lines 2769--2769. Pointer outside the counted span, invoked once for context; the load-bearing argument of the counted statement is the author's proof reproduced in full in the proof span.\par\medskip
\begin{definition}
\label{definition-affine-formal-algebraic-space}
Let $S$ be a scheme. We say a sheaf $X$ on $(\Sch/S)_{fppf}$ is an
{\it affine formal algebraic space} if there exist
\begin{enumerate}
\item a directed set $\Lambda$,
\item a system $(X_\lambda, f_{\lambda \mu})$ over $\Lambda$
in $(\Sch/S)_{fppf}$ where
\begin{enumerate}
\item each $X_\lambda$ is affine,
\item each $f_{\lambda \mu} : X_\lambda \to X_\mu$ is a thickening,
\end{enumerate}
\end{enumerate}
such that
$$
X \cong \colim_{\lambda \in \Lambda} X_\lambda
$$
as fppf sheaves and $X$ satisfies a set theoretic condition
(see Remark \ref{remark-set-theoretic}). A
{\it morphism of affine formal algebraic spaces}
over $S$ is a map of sheaves.
\end{definition}

\begin{lemma}
\label{lemma-structure-sheaf}
Every formal algebraic space has a structure sheaf.
\end{lemma}

\begin{proof}
Let $S$ be a scheme. Let $X$ be a formal algebraic space over $S$.
By (\ref{equation-etale-topos}) it suffices to construct $\mathcal{O}_X$
as a sheaf of topological rings on $X_{affine, \etale}$.
Denote $\mathcal{C}$ the category whose objects
are morphisms $\varphi : U \to X$ of formal algebraic spaces such that
$U$ is an affine formal algebraic space and
$\varphi$ is representable by algebraic spaces and \'etale.
By Lemma \ref{lemma-identify-affine-etale}
the functor $U \mapsto U_{red}$ is an equivalence of categories
$\mathcal{C} \to X_{affine, \etale}$. Hence by the rule given
above the lemma, we already have $\mathcal{O}_X$ as a presheaf of
topological rings on $X_{affine, \etale}$. Thus it suffices to check
the sheaf condition.

\medskip\noindent
By definition of $X_{affine, \etale}$ a covering corresponds to a finite
family $\{g_i : U_i \to U\}_{i = 1, \ldots, n}$ of morphisms of $\mathcal{C}$
such that $\{U_{i, red} \to U_{red}\}$ is an \'etale covering.
The morphisms $g_i$ are representably by algebraic spaces
(Lemma \ref{lemma-permanence-representable}) hence affine
(Lemma \ref{lemma-affine-representable-by-algebraic-spaces}).
Then $g_i$ is \'etale (follows formally from
Properties of Spaces, Lemma \ref{spaces-properties-lemma-etale-permanence}
as $U_i$ and $U$ are \'etale over $X$ in the sense of
Bootstrap, Section \ref{bootstrap-section-representable-by-spaces-properties}).
Finally, write $U = \colim U_\lambda$ as in
Definition \ref{definition-affine-formal-algebraic-space}.

\medskip\noindent
With these preparations out of the way, we can prove the sheaf
property as follows. For each $\lambda$ we set
$U_{i, \lambda} = U_i \times_U U_\lambda$ and
$U_{ij, \lambda} = (U_i \times_U U_j) \times_U U_\lambda$.
By the above, these are affine schemes, $\{U_{i, \lambda} \to U_\lambda\}$
is an \'etale covering, and
$U_{ij, \lambda} = U_{i, \lambda} \times_{U_\lambda} U_{j, \lambda}$.
Also we have $U_i = \colim U_{i, \lambda}$ and
$U_i \times_U U_j = \colim U_{ij, \lambda}$.
For each $\lambda$ we have an exact sequence
$$
0 \to
\Gamma(U_\lambda, \mathcal{O}_{U_\lambda}) \to
\prod\nolimits_i \Gamma(U_{i, \lambda}, \mathcal{O}_{U_{i, \lambda}}) \to
\prod\nolimits_{i, j} \Gamma(U_{ij, \lambda}, \mathcal{O}_{U_{ij, \lambda}})
$$
as we have the sheaf condition for the structure sheaf
on $U_\lambda$ and the \'etale topology
(see \'Etale Cohomology, Proposition
\ref{etale-cohomology-proposition-quasi-coherent-sheaf-fpqc}).
Since limits commute with limits, the inverse limit of these
exact sequences is an exact sequence
$$
0 \to
\lim \Gamma(U_\lambda, \mathcal{O}_{U_\lambda}) \to
\prod\nolimits_i \lim \Gamma(U_{i, \lambda}, \mathcal{O}_{U_{i, \lambda}}) \to
\prod\nolimits_{i, j} \lim
\Gamma(U_{ij, \lambda}, \mathcal{O}_{U_{ij, \lambda}})
$$
This means that
$$
0 \to
\mathcal{O}_X(U_{red}) \to
\prod\nolimits_i \mathcal{O}_X(U_{i, red}) \to
\prod\nolimits_{i, j} \mathcal{O}_X((U_i \times_U U_j)_{red})
$$
is exact as a sequence of abelian groups. We still have to
show that the limit topology on $\mathcal{O}_X(U_{red})$ is induced
from the limit topology on
$\prod_{i = 1, \ldots, n} \mathcal{O}_X(U_{i, red})$.
The kernels of the maps
$\mathcal{O}_X(U_{red}) \to \mathcal{O}_{U_\lambda}(U_\lambda)$
form a fundamental system of open submodules. Since these are the inverse
images of the kernels of the maps
$\prod_i \mathcal{O}_X(U_{i, red}) \to
\prod_i \mathcal{O}_{U_{i, \lambda}}(U_{i, \lambda})$
and these kernels form a fundamental system of open submodules
of $\mathcal{O}_X(U_{red}) \to \mathcal{O}_{U_\lambda}(U_\lambda)$
we win.
\end{proof}

\end{document}
